\documentclass{umk-matematika}

\begin{document}

\maketitlepage
    {Практическая работа №5}
    {Аксиомы стереометрии. Сечения}
    {}

\section{Фундамент (1--4 балла)}

\textbf{Задача 1 (1 балл).} Сформулируйте аксиому 1 стереометрии.

\textbf{Задача 2 (2 балла).} Сформулируйте аксиому 2 стереометрии.

\textbf{Задача 3 (3 балла).} Сформулируйте аксиому 3 стереометрии.

\textbf{Задача 4 (4 балла).} Сколько плоскостей можно провести через две пересекающиеся прямые?

\section{Стены (5--7 баллов)}

\textbf{Задача 5 (5 баллов).} Могут ли три точки одновременно принадлежать двум разным плоскостям?

\textbf{Задача 6 (6 баллов).} Даны четыре точки, не лежащие в одной плоскости. Сколько различных плоскостей можно провести через различные тройки этих точек?

\textbf{Задача 7 (7 баллов).} Что такое сечение многогранника? Какие многоугольники могут быть сечениями куба?

\section{Кровля (8--10 баллов)}

\textbf{Задача 8 (8 баллов).} Постройте сечение куба $ABCDA_1B_1C_1D_1$ плоскостью, проходящей через середины рёбер $AB$, $BC$ и вершину $D_1$. Какой многоугольник получится?

\textbf{Задача 9 (9 баллов).} Какие многоугольники могут быть сечениями тетраэдра?

\textbf{Задача 10 (10 баллов).} Строительный блок имеет форму прямоугольного параллелепипеда. Его распилили плоскостью, проходящей через диагонали двух противоположных граней. Какая фигура получилась в сечении?

\newpage
\section*{Ответы}

\begin{enumerate}
  \item аксиома 1.
  \item аксиома 2.
  \item аксиома 3.
  \item одну.
  \item да, если точки лежат на одной прямой; нет, если не лежат.
  \item 4 плоскости.
  \item треугольник, четырёхугольник, пятиугольник, шестиугольник.
  \item пятиугольник.
  \item треугольник и четырёхугольник.
  \item прямоугольник.
\end{enumerate}

\end{document}